\section{Metodología}

\subsection{Materiales y Herramientas}

\begin{itemize}
    \item 
\end{itemize}



\subsection{Procedimiento general}

\subsubsection{Plotter} ---
El sistema del plotter tiene como objetivo controlar de manera coordinada dos motores paso a paso y una válvula solenoide utilizando el microcontrolador ATmega328P. A través de la interfaz serial (UART), el usuario puede seleccionar distintas figuras predeterminadas —como triángulos, cruces o círculos— para su trazado automático. El microcontrolador envía las órdenes de dirección y activación al sistema de potencia, asegurando que los movimientos se mantengan dentro de los límites definidos por los sensores de fin de carrera. De esta forma, el dispositivo permite realizar dibujos controlados en los ejes X e Y, integrando comunicación serial, control de motores y manejo de periféricos en una aplicación práctica de control embebido.

\vspace{1em}

El programa implementa el control completo de un \textit{plotter} didáctico de dos ejes (X e Y) utilizando motores paso a paso y un solenoide que sube o baja el lápiz para dibujar figuras a partir de trayectorias predefinidas almacenadas en memoria de programa (\texttt{path\_data.h}). Emplea interrupciones externas para los finales de carrera (PD2 y PD3), junto con un sistema de antirrebote basado en el temporizador \texttt{Timer0}, que asegura lecturas estables. La función principal inicializa los periféricos y ejecuta de forma secuencial distintas trayectorias (murciélago, círculo, flor, cruz y triángulo), repitiendo el ciclo de dibujo. Cada trayectoria se realiza mediante la función \texttt{execute\_path()}, que mueve los motores de acuerdo con los desplazamientos definidos, controlando dirección y cantidad de pasos en cada eje, mientras el solenoide indica cuándo dibujar. Si un límite se activa, el sistema detiene el movimiento mediante \texttt{emergency\_stop()}, garantizando seguridad. En conjunto, el código conforma un sistema de control secuencial de dibujo automático con soporte para trayectorias reflejadas y detección de límites mecánicos.


\vspace{1em}
\subsubsection{Temperatura} ---
Este módulo mide la temperatura ambiente mediante un sensor LM35 y ajusta automáticamente el funcionamiento de un calefactor y un ventilador según el rango térmico detectado. El ATmega328P ejecuta lecturas periódicas a través del ADC y determina las acciones de control utilizando salidas digitales y señales PWM. El sistema comunica los valores y estados por UART, e incluye un menú que permite modificar el punto medio de temperatura. Además, se prevé una simulación y una implementación física del sistema, complementadas con una gráfica de evolución de temperatura generada en Python o MATLAB.

\vspace{1em}

El programa implementa un sistema de \textbf{control automático de temperatura} utilizando un microcontrolador ATmega328P, un sensor \textit{LM35} para la medición analógica de temperatura y un \textit{ventilador controlado por PWM}. La lectura del sensor se realiza mediante el módulo ADC, y el valor obtenido se convierte a grados Celsius para ser comparado con los umbrales configurables \( t_{C2} \) (mínimo) y \( t_{C3} \) (máximo). Si la temperatura medida excede el rango establecido, el ventilador se activa o desactiva automáticamente. La comunicación con el usuario se realiza a través de la interfaz \textit{UART}, permitiendo visualizar la temperatura, modificar los límites y controlar el sistema mediante comandos como \texttt{Encender}. El programa utiliza interrupciones por temporizador para realizar mediciones periódicas y \textit{buffers} circulares para gestionar la comunicación serial de forma no bloqueante. En conjunto, el sistema presenta un funcionamiento modular, estable y adaptable a diferentes condiciones térmicas.


\vspace{1em}
\subsubsection{Motor} ---
En este sistema se comparan las lecturas de dos potenciómetros, uno como referencia y otro acoplado al eje de un motor, con el fin de igualar ambas señales. El microcontrolador utiliza el conversor ADC para leer los valores y genera una señal PWM que ajusta la velocidad y el sentido de giro del motor hasta lograr la coincidencia. Durante el proceso, se transmiten por UART los valores de referencia, posición actual, nivel de PWM y sentido de giro, permitiendo un seguimiento en tiempo real. El sistema se implementa tanto en simulación como físicamente, complementado con una visualización gráfica de los resultados.

\vspace{1em}

El programa implementa el \textbf{control de posición de un motor de corriente continua} con retroalimentación analógica, utilizando el microcontrolador \textit{ATmega328P}. El sistema compara en tiempo real dos señales analógicas leídas por el ADC: una de referencia (posición deseada) y otra proveniente del potenciómetro acoplado al motor (posición actual). En función de la diferencia o \textit{error}, el microcontrolador ajusta la dirección de giro mediante los pines \texttt{IN1} e \texttt{IN2}, y genera una señal PWM proporcional en \texttt{OCR1A} para regular la velocidad. Se establece una \textit{zona muerta} para evitar oscilaciones y se limitan los valores mínimos y máximos de PWM. Además, el sistema transmite por \textit{USART} los valores de referencia, posición real y salida PWM, junto con la dirección de movimiento, permitiendo el monitoreo en tiempo real desde un ordenador. El procesamiento principal se ejecuta dentro de un bucle que llama periódicamente a la función \texttt{usart\_task()}, mientras las interrupciones gestionan la comunicación serial mediante \textit{buffers} circulares no bloqueantes.


\vspace{1em}
\subsubsection{Matriz} ---
El sistema emplea un joystick analógico para desplazar un LED activo dentro de una matriz RGB, utilizando el ADC del ATmega328P para detectar los movimientos sobre los ejes X e Y. Los límites de desplazamiento están restringidos al área visible de la matriz, y al presionar el botón del joystick se cambia aleatoriamente el color del LED activo. Se incorpora una zona muerta para evitar lecturas erráticas por ruido analógico, garantizando un control preciso y fluido. Esta parte se desarrolla únicamente de forma física en el laboratorio.

\vspace{1em}

El programa implementa un sistema de \textbf{control interactivo de una matriz de LEDs WS2812} (64 píxeles dispuestos en una cuadrícula de 8×8) utilizando el microcontrolador \textit{ATmega328P}, un \textit{joystick analógico} y un \textit{botón con antirrebote}. El joystick controla la posición de un píxel encendido dentro de la matriz (variables \texttt{x\_pos} y \texttt{y\_pos}), mientras que el botón genera colores aleatorios mediante una rutina de interrupción con \textit{debounce}. La comunicación serial \textit{USART} (reutilizada del laboratorio 2) permite monitorear el sistema o enviar datos de depuración. Para generar las señales temporizadas requeridas por los LEDs WS2812, se emplean rutinas en ensamblador inline con \texttt{asm volatile}, asegurando los ciclos exactos por bit. Además, el código utiliza interrupciones por temporizador para contar milisegundos (\texttt{millis\_counter}) y gestionar el antirrebote por software. En conjunto, el sistema integra control analógico, comunicación serial y temporización precisa para crear una interfaz visual interactiva sobre una matriz RGB.


\vspace{1em}
\subsubsection{Cerradura} ---
La cerradura electrónica implementa un sistema de autenticación mediante tarjetas RFID RC522 controladas por el microcontrolador ATmega328P. El identificador único (ID) de la tarjeta autorizada se almacena en la memoria EEPROM para su verificación permanente. La interfaz incluye una pantalla LCD I2C, LEDs indicadores y comunicación UART para mensajes informativos y supervisión externa. El sistema permite registrar, actualizar y eliminar tarjetas de acceso, combinando almacenamiento no volátil, comunicación serie y periféricos de entrada/salida en un diseño de seguridad embebido. Este módulo se desarrolla tanto en simulación como en implementación física.

El programa implementa un sistema de \textbf{control de acceso mediante RFID} con visualización en pantalla LCD e interfaz serial, utilizando el microcontrolador \textit{ATmega328P} y el módulo \textit{MFRC522} para la lectura de tarjetas. El sistema cuenta con tres modos principales: verificación de acceso, registro de una nueva tarjeta y borrado del UID almacenado. Los identificadores se guardan de forma persistente en la memoria \textit{EEPROM}, y se comparan con los leídos en tiempo real. La interacción del usuario se realiza tanto por medio de la comunicación \textit{UART} como mediante botones físicos conectados al puerto C. La información de estado se muestra en una pantalla LCD controlada por interfaz \textit{I2C} (expansor PCF8574), mientras que un buzzer emite alertas acústicas en casos de acceso denegado o eventos relevantes. Internamente, el sistema utiliza comunicación \textit{SPI} para el módulo RFID, \textit{I2C} para el LCD y \textit{USART} con buffers circulares para comunicación no bloqueante. En conjunto, el código integra múltiples periféricos y subsistemas (SPI, I2C, UART, EEPROM e interrupciones) para construir un sistema embebido completo y robusto de autenticación electrónica.


\vspace{1em}
\subsubsection{Fútbol} ---
Esta consigna consiste en el diseño y programación de un robot móvil basado en el ATmega328P, capaz de disputar partidos de mini fútbol sobre un campo reducido. El robot integra sensores, actuadores y sistemas de comunicación inalámbrica (Bluetooth o RF) para detectar y conducir el balón hacia la portería rival. Se evalúan la autonomía, el control, la integración mecánica y la creatividad del diseño. Este proyecto se realiza únicamente de forma física, siendo el cierre práctico y competitivo del laboratorio.
